% Two-panel optimization analogy for Section 1.2.2.
% Requires \usepackage{tikz}; \usetikzlibrary{arrows.meta,positioning,calc}
% Insert using % Two-panel optimization analogy for Section 1.2.2.
% Requires \usepackage{tikz}; \usetikzlibrary{arrows.meta,positioning,calc}
% Insert using % Two-panel optimization analogy for Section 1.2.2.
% Requires \usepackage{tikz}; \usetikzlibrary{arrows.meta,positioning,calc}
% Insert using % Two-panel optimization analogy for Section 1.2.2.
% Requires \usepackage{tikz}; \usetikzlibrary{arrows.meta,positioning,calc}
% Insert using \input{optimization_figure.tex} after the two-paragraph
% discussion beginning "Training as an Optimization Problem".
\begin{figure}[tbp]
\centering
\resizebox{\textwidth}{!}{%
\begin{tikzpicture}[font=\sffamily\footnotesize,>=Latex,
alt={Two examples of optimization. On the left, changing the visiting order produces a shorter traveling-salesman route. On the right, adjusting neural-network weights reduces loss along an illustrative direction in parameter space.}]
  % Panel (a): a route and a nearby shorter route, drawn with the same cities.
  \begin{scope}[shift={(0,0)}]
    \node[anchor=west,font=\sffamily\bfseries] at (0,3.65)
      {(a) Traveling salesman: improve a route};
    \node[font=\scriptsize] at (1.65,3.20) {Initial route};
    \node[font=\scriptsize] at (5.05,3.20) {Shorter route};
    % The original route has a crossing; the new route removes it
    % by swapping the order of two stops.
    \begin{scope}[shift={(0,.35)}]
      \coordinate (A) at (.50,.25);
      \coordinate (B) at (2.48,.48);
      \coordinate (C) at (2.25,2.20);
      \coordinate (D) at (.93,2.49);
      \coordinate (E) at (.12,1.37);
      \draw[orange!85!black,line width=1.15pt] (A)--(B)--(D)--(C)--(E)--cycle;
      \foreach \p/\name in {A/A,B/B,C/C,D/D,E/E} {
        \draw[fill=white,draw=blue!70!black,line width=.7pt] (\p) circle (2.8pt);
        \node[font=\tiny,anchor=west] at ($ (\p)+(.09,.02) $) {\name};
      }
    \end{scope}
    \begin{scope}[shift={(3.4,.35)}]
      \coordinate (A) at (.50,.25);
      \coordinate (B) at (2.48,.48);
      \coordinate (C) at (2.25,2.20);
      \coordinate (D) at (.93,2.49);
      \coordinate (E) at (.12,1.37);
      \draw[teal!70!black,line width=1.15pt] (A)--(B)--(C)--(D)--(E)--cycle;
      \foreach \p/\name in {A/A,B/B,C/C,D/D,E/E} {
        \draw[fill=white,draw=blue!70!black,line width=.7pt] (\p) circle (2.8pt);
        \node[font=\tiny,anchor=west] at ($ (\p)+(.09,.02) $) {\name};
      }
    \end{scope}
    \draw[->,gray!70!black,line width=.9pt] (2.91,1.65)--(3.27,1.65);
    \node[align=center,font=\scriptsize,text=black!75] at (3.3,-.03)
       {Change the visiting order; compare total distances.};
  \end{scope}

  % Panel (b): illustrative cross-section of a continuously varying loss.
  \begin{scope}[shift={(7.65,0)}]
    \node[anchor=west,font=\sffamily\bfseries] at (0,3.65)
      {(b) Neural-network training: lower the loss};
    \draw[->,gray!80!black] (.48,.46)--(6.40,.46);
    \draw[->,gray!80!black] (.48,.46)--(.48,3.16);
    \node[anchor=west,font=\scriptsize] at (.53,3.08) {Loss (prediction error)};
    \node[font=\scriptsize] at (3.55,.07) {A model-weight direction (schematic)};
    % Same curve for all points: y = .62 + .17*(x-3.25)^2
    \draw[blue!75!black,thick,domain=.72:6.20,samples=75,smooth]
      plot (\x,{.63+.17*(\x-3.25)*(\x-3.25)});
    % Three downhill steps (right to left), plus final near the minimum.
    \coordinate (p1) at (5.75,1.6925);
    \coordinate (p2) at (4.94,1.1156);
    \coordinate (p3) at (4.18,.7770);
    \coordinate (p4) at (3.62,.6533);
    \draw[->,orange!90!black,line width=1pt] (p1) -- (p2);
    \draw[->,orange!90!black,line width=1pt] (p2) -- (p3);
    \draw[->,orange!90!black,line width=1pt] (p3) -- (p4);
    \foreach \p in {p1,p2,p3,p4} {
      \draw[fill=orange!85!black,draw=white,line width=.5pt] (\p) circle (2.6pt);
    }
    \node[font=\scriptsize,anchor=west] at (5.17,2.20) {Start};
  \end{scope}
\end{tikzpicture}%
}
\caption{Two kinds of optimization. (a) In the traveling salesman problem,
changing the order of visits can shorten a route through the same cities.
The alternatives are discrete routes. (b) In neural-network training,
small adjustments to numerical weights can reduce a measure of prediction
error (the \emph{loss}). The curve represents one illustrative direction
through a much higher-dimensional loss landscape; gradient-based training
does not guarantee a global minimum.}
\label{fig:optimization}
\end{figure}
 after the two-paragraph
% discussion beginning "Training as an Optimization Problem".
\begin{figure}[tbp]
\centering
\resizebox{\textwidth}{!}{%
\begin{tikzpicture}[font=\sffamily\footnotesize,>=Latex,
alt={Two examples of optimization. On the left, changing the visiting order produces a shorter traveling-salesman route. On the right, adjusting neural-network weights reduces loss along an illustrative direction in parameter space.}]
  % Panel (a): a route and a nearby shorter route, drawn with the same cities.
  \begin{scope}[shift={(0,0)}]
    \node[anchor=west,font=\sffamily\bfseries] at (0,3.65)
      {(a) Traveling salesman: improve a route};
    \node[font=\scriptsize] at (1.65,3.20) {Initial route};
    \node[font=\scriptsize] at (5.05,3.20) {Shorter route};
    % The original route has a crossing; the new route removes it
    % by swapping the order of two stops.
    \begin{scope}[shift={(0,.35)}]
      \coordinate (A) at (.50,.25);
      \coordinate (B) at (2.48,.48);
      \coordinate (C) at (2.25,2.20);
      \coordinate (D) at (.93,2.49);
      \coordinate (E) at (.12,1.37);
      \draw[orange!85!black,line width=1.15pt] (A)--(B)--(D)--(C)--(E)--cycle;
      \foreach \p/\name in {A/A,B/B,C/C,D/D,E/E} {
        \draw[fill=white,draw=blue!70!black,line width=.7pt] (\p) circle (2.8pt);
        \node[font=\tiny,anchor=west] at ($ (\p)+(.09,.02) $) {\name};
      }
    \end{scope}
    \begin{scope}[shift={(3.4,.35)}]
      \coordinate (A) at (.50,.25);
      \coordinate (B) at (2.48,.48);
      \coordinate (C) at (2.25,2.20);
      \coordinate (D) at (.93,2.49);
      \coordinate (E) at (.12,1.37);
      \draw[teal!70!black,line width=1.15pt] (A)--(B)--(C)--(D)--(E)--cycle;
      \foreach \p/\name in {A/A,B/B,C/C,D/D,E/E} {
        \draw[fill=white,draw=blue!70!black,line width=.7pt] (\p) circle (2.8pt);
        \node[font=\tiny,anchor=west] at ($ (\p)+(.09,.02) $) {\name};
      }
    \end{scope}
    \draw[->,gray!70!black,line width=.9pt] (2.91,1.65)--(3.27,1.65);
    \node[align=center,font=\scriptsize,text=black!75] at (3.3,-.03)
       {Change the visiting order; compare total distances.};
  \end{scope}

  % Panel (b): illustrative cross-section of a continuously varying loss.
  \begin{scope}[shift={(7.65,0)}]
    \node[anchor=west,font=\sffamily\bfseries] at (0,3.65)
      {(b) Neural-network training: lower the loss};
    \draw[->,gray!80!black] (.48,.46)--(6.40,.46);
    \draw[->,gray!80!black] (.48,.46)--(.48,3.16);
    \node[anchor=west,font=\scriptsize] at (.53,3.08) {Loss (prediction error)};
    \node[font=\scriptsize] at (3.55,.07) {A model-weight direction (schematic)};
    % Same curve for all points: y = .62 + .17*(x-3.25)^2
    \draw[blue!75!black,thick,domain=.72:6.20,samples=75,smooth]
      plot (\x,{.63+.17*(\x-3.25)*(\x-3.25)});
    % Three downhill steps (right to left), plus final near the minimum.
    \coordinate (p1) at (5.75,1.6925);
    \coordinate (p2) at (4.94,1.1156);
    \coordinate (p3) at (4.18,.7770);
    \coordinate (p4) at (3.62,.6533);
    \draw[->,orange!90!black,line width=1pt] (p1) -- (p2);
    \draw[->,orange!90!black,line width=1pt] (p2) -- (p3);
    \draw[->,orange!90!black,line width=1pt] (p3) -- (p4);
    \foreach \p in {p1,p2,p3,p4} {
      \draw[fill=orange!85!black,draw=white,line width=.5pt] (\p) circle (2.6pt);
    }
    \node[font=\scriptsize,anchor=west] at (5.17,2.20) {Start};
  \end{scope}
\end{tikzpicture}%
}
\caption{Two kinds of optimization. (a) In the traveling salesman problem,
changing the order of visits can shorten a route through the same cities.
The alternatives are discrete routes. (b) In neural-network training,
small adjustments to numerical weights can reduce a measure of prediction
error (the \emph{loss}). The curve represents one illustrative direction
through a much higher-dimensional loss landscape; gradient-based training
does not guarantee a global minimum.}
\label{fig:optimization}
\end{figure}
 after the two-paragraph
% discussion beginning "Training as an Optimization Problem".
\begin{figure}[tbp]
\centering
\resizebox{\textwidth}{!}{%
\begin{tikzpicture}[font=\sffamily\footnotesize,>=Latex,
alt={Two examples of optimization. On the left, changing the visiting order produces a shorter traveling-salesman route. On the right, adjusting neural-network weights reduces loss along an illustrative direction in parameter space.}]
  % Panel (a): a route and a nearby shorter route, drawn with the same cities.
  \begin{scope}[shift={(0,0)}]
    \node[anchor=west,font=\sffamily\bfseries] at (0,3.65)
      {(a) Traveling salesman: improve a route};
    \node[font=\scriptsize] at (1.65,3.20) {Initial route};
    \node[font=\scriptsize] at (5.05,3.20) {Shorter route};
    % The original route has a crossing; the new route removes it
    % by swapping the order of two stops.
    \begin{scope}[shift={(0,.35)}]
      \coordinate (A) at (.50,.25);
      \coordinate (B) at (2.48,.48);
      \coordinate (C) at (2.25,2.20);
      \coordinate (D) at (.93,2.49);
      \coordinate (E) at (.12,1.37);
      \draw[orange!85!black,line width=1.15pt] (A)--(B)--(D)--(C)--(E)--cycle;
      \foreach \p/\name in {A/A,B/B,C/C,D/D,E/E} {
        \draw[fill=white,draw=blue!70!black,line width=.7pt] (\p) circle (2.8pt);
        \node[font=\tiny,anchor=west] at ($ (\p)+(.09,.02) $) {\name};
      }
    \end{scope}
    \begin{scope}[shift={(3.4,.35)}]
      \coordinate (A) at (.50,.25);
      \coordinate (B) at (2.48,.48);
      \coordinate (C) at (2.25,2.20);
      \coordinate (D) at (.93,2.49);
      \coordinate (E) at (.12,1.37);
      \draw[teal!70!black,line width=1.15pt] (A)--(B)--(C)--(D)--(E)--cycle;
      \foreach \p/\name in {A/A,B/B,C/C,D/D,E/E} {
        \draw[fill=white,draw=blue!70!black,line width=.7pt] (\p) circle (2.8pt);
        \node[font=\tiny,anchor=west] at ($ (\p)+(.09,.02) $) {\name};
      }
    \end{scope}
    \draw[->,gray!70!black,line width=.9pt] (2.91,1.65)--(3.27,1.65);
    \node[align=center,font=\scriptsize,text=black!75] at (3.3,-.03)
       {Change the visiting order; compare total distances.};
  \end{scope}

  % Panel (b): illustrative cross-section of a continuously varying loss.
  \begin{scope}[shift={(7.65,0)}]
    \node[anchor=west,font=\sffamily\bfseries] at (0,3.65)
      {(b) Neural-network training: lower the loss};
    \draw[->,gray!80!black] (.48,.46)--(6.40,.46);
    \draw[->,gray!80!black] (.48,.46)--(.48,3.16);
    \node[anchor=west,font=\scriptsize] at (.53,3.08) {Loss (prediction error)};
    \node[font=\scriptsize] at (3.55,.07) {A model-weight direction (schematic)};
    % Same curve for all points: y = .62 + .17*(x-3.25)^2
    \draw[blue!75!black,thick,domain=.72:6.20,samples=75,smooth]
      plot (\x,{.63+.17*(\x-3.25)*(\x-3.25)});
    % Three downhill steps (right to left), plus final near the minimum.
    \coordinate (p1) at (5.75,1.6925);
    \coordinate (p2) at (4.94,1.1156);
    \coordinate (p3) at (4.18,.7770);
    \coordinate (p4) at (3.62,.6533);
    \draw[->,orange!90!black,line width=1pt] (p1) -- (p2);
    \draw[->,orange!90!black,line width=1pt] (p2) -- (p3);
    \draw[->,orange!90!black,line width=1pt] (p3) -- (p4);
    \foreach \p in {p1,p2,p3,p4} {
      \draw[fill=orange!85!black,draw=white,line width=.5pt] (\p) circle (2.6pt);
    }
    \node[font=\scriptsize,anchor=west] at (5.17,2.20) {Start};
  \end{scope}
\end{tikzpicture}%
}
\caption{Two kinds of optimization. (a) In the traveling salesman problem,
changing the order of visits can shorten a route through the same cities.
The alternatives are discrete routes. (b) In neural-network training,
small adjustments to numerical weights can reduce a measure of prediction
error (the \emph{loss}). The curve represents one illustrative direction
through a much higher-dimensional loss landscape; gradient-based training
does not guarantee a global minimum.}
\label{fig:optimization}
\end{figure}
 after the two-paragraph
% discussion beginning "Training as an Optimization Problem".
\begin{figure}[tbp]
\centering
\resizebox{\textwidth}{!}{%
\begin{tikzpicture}[font=\sffamily\footnotesize,>=Latex,
alt={Two examples of optimization. On the left, changing the visiting order produces a shorter traveling-salesman route. On the right, adjusting neural-network weights reduces loss along an illustrative direction in parameter space.}]
  % Panel (a): a route and a nearby shorter route, drawn with the same cities.
  \begin{scope}[shift={(0,0)}]
    \node[anchor=west,font=\sffamily\bfseries] at (0,3.65)
      {(a) Traveling salesman: improve a route};
    \node[font=\scriptsize] at (1.65,3.20) {Initial route};
    \node[font=\scriptsize] at (5.05,3.20) {Shorter route};
    % The original route has a crossing; the new route removes it
    % by swapping the order of two stops.
    \begin{scope}[shift={(0,.35)}]
      \coordinate (A) at (.50,.25);
      \coordinate (B) at (2.48,.48);
      \coordinate (C) at (2.25,2.20);
      \coordinate (D) at (.93,2.49);
      \coordinate (E) at (.12,1.37);
      \draw[orange!85!black,line width=1.15pt] (A)--(B)--(D)--(C)--(E)--cycle;
      \foreach \p/\name in {A/A,B/B,C/C,D/D,E/E} {
        \draw[fill=white,draw=blue!70!black,line width=.7pt] (\p) circle (2.8pt);
        \node[font=\tiny,anchor=west] at ($ (\p)+(.09,.02) $) {\name};
      }
    \end{scope}
    \begin{scope}[shift={(3.4,.35)}]
      \coordinate (A) at (.50,.25);
      \coordinate (B) at (2.48,.48);
      \coordinate (C) at (2.25,2.20);
      \coordinate (D) at (.93,2.49);
      \coordinate (E) at (.12,1.37);
      \draw[teal!70!black,line width=1.15pt] (A)--(B)--(C)--(D)--(E)--cycle;
      \foreach \p/\name in {A/A,B/B,C/C,D/D,E/E} {
        \draw[fill=white,draw=blue!70!black,line width=.7pt] (\p) circle (2.8pt);
        \node[font=\tiny,anchor=west] at ($ (\p)+(.09,.02) $) {\name};
      }
    \end{scope}
    \draw[->,gray!70!black,line width=.9pt] (2.91,1.65)--(3.27,1.65);
    \node[align=center,font=\scriptsize,text=black!75] at (3.3,-.03)
       {Change the visiting order; compare total distances.};
  \end{scope}

  % Panel (b): illustrative cross-section of a continuously varying loss.
  \begin{scope}[shift={(7.65,0)}]
    \node[anchor=west,font=\sffamily\bfseries] at (0,3.65)
      {(b) Neural-network training: lower the loss};
    \draw[->,gray!80!black] (.48,.46)--(6.40,.46);
    \draw[->,gray!80!black] (.48,.46)--(.48,3.16);
    \node[anchor=west,font=\scriptsize] at (.53,3.08) {Loss (prediction error)};
    \node[font=\scriptsize] at (3.55,.07) {A model-weight direction (schematic)};
    % Same curve for all points: y = .62 + .17*(x-3.25)^2
    \draw[blue!75!black,thick,domain=.72:6.20,samples=75,smooth]
      plot (\x,{.63+.17*(\x-3.25)*(\x-3.25)});
    % Three downhill steps (right to left), plus final near the minimum.
    \coordinate (p1) at (5.75,1.6925);
    \coordinate (p2) at (4.94,1.1156);
    \coordinate (p3) at (4.18,.7770);
    \coordinate (p4) at (3.62,.6533);
    \draw[->,orange!90!black,line width=1pt] (p1) -- (p2);
    \draw[->,orange!90!black,line width=1pt] (p2) -- (p3);
    \draw[->,orange!90!black,line width=1pt] (p3) -- (p4);
    \foreach \p in {p1,p2,p3,p4} {
      \draw[fill=orange!85!black,draw=white,line width=.5pt] (\p) circle (2.6pt);
    }
    \node[font=\scriptsize,anchor=west] at (5.17,2.20) {Start};
  \end{scope}
\end{tikzpicture}%
}
\caption{Two kinds of optimization. (a) In the traveling salesman problem,
changing the order of visits can shorten a route through the same cities.
The alternatives are discrete routes. (b) In neural-network training,
small adjustments to numerical weights can reduce a measure of prediction
error (the \emph{loss}). The curve represents one illustrative direction
through a much higher-dimensional loss landscape; gradient-based training
does not guarantee a global minimum.}
\label{fig:optimization}
\end{figure}
